% P15.31: Durrett's general central-limit presentation.
\chapter{Central Limit Theorems}
\label{sec:p15-central-limit-theorems}

A central limit theorem identifies the distribution of a centered sum
at its fluctuation scale. Independence within each row, control of its
total variance, and exclusion of large individual contributions are
different requirements. None of them follows merely from a law of large
numbers.

This chapter develops the source's characteristic-function and
method-of-moments proofs, its applications and exercises, and its
quantitative normal approximation
\cite[Section~3.4, pp.~144--161]{DurrettPTEAuthor2019}.
The source version is the author Version 5 dated 11 January 2019.

\begin{hbdefinition}{def:p15-clt-notation}{Normal limits and row normalization}
All random variables are finite, real, and measurable. Write
\[
 \phi_0(x)=(2\pi)^{-1/2}e^{-x^2/2},\qquad
 \Phi(x)=\int_{-\infty}^x\phi_0(u)\,du,
\]
and let \(\chi\) denote a variable with this standard normal law.
The notation \(Z_n\Rightarrow Z\) means weak convergence of laws.
The law \(N(0,\sigma^2)\) means that of \(\sigma\chi\);
at \(\sigma=0\) it is point mass at zero.
Whenever a single sequence \(X_1,X_2,\ldots\) is under
discussion, its partial sums are \(S_0=0\) and
\(S_n=\sum_{j=1}^nX_j\).

A triangular array has entries \(X_{n,j}\), \(1\leq j\leq k_n\),
where each \(k_n\) is a positive finite integer. Row independence
means mutual independence within each fixed row; no independence
between rows is required. Put \(S_n=\sum_{j=1}^{k_n}X_{n,j}\).
For centered square-integrable entries set
\[
 v_{n,j}=\mathbb EX_{n,j}^2,\qquad
 v_n=\sum_{j=1}^{k_n}v_{n,j},\qquad
 L_n(\varepsilon)=
 \sum_{j=1}^{k_n}\mathbb E(X_{n,j}^2;
                                  |X_{n,j}|>\varepsilon).
\]
The notation \(\mathbb E(W;A)\) means
\(\mathbb E(W\mathbf1_A)\). All strict cutoff inequalities in
\(L_n\) are part of its definition.
For positive deterministic scales \(c_n\), write
\(R_n=o_{\mathbb P}(c_n)\) when \(R_n/c_n\to0\) in probability,
and \(R_n=O_{\mathbb P}(c_n)\) when the laws of \(R_n/c_n\)
are tight. A family is tight when for every \(\varepsilon>0\)
some finite \(K\) bounds all its probabilities of \(\{|Z|>K\}\)
by \(\varepsilon\). The same notation applies to a real parameter
tending to infinity. For deterministic quantities, \(u_n\sim v_n\)
means \(u_n/v_n\to1\), where \(v_n\ne0\) eventually.
\end{hbdefinition}

\section{Independent identically distributed summands}
\label{sec:p15-clt-iid}

\begin{hblemma}{thm:p15-clt-complex-products}{proved=annex:p15.clt-complex-products}{The source's complex-product argument}
For \(n\geq1\) and complex numbers \(z_j,w_j\) with
\(|z_j|,|w_j|\leq\theta\), where \(\theta>0\),
\[
 \left|\prod_{j=1}^nz_j-\prod_{j=1}^nw_j\right|
 \leq\theta^{n-1}\sum_{j=1}^n|z_j-w_j|.
\]
If \(\theta=0\), every factor is zero and the product difference
is zero; this is the degenerate case of the same estimate.
For \(|b|\leq1\), \(|e^b-1-b|\leq|b|^2\).
Consequently, \(c_n\to c\) in \(\mathbb C\) implies
\((1+c_n/n)^n\to e^c\).
These are Lemmas 3.4.3--3.4.4 and Theorem 3.4.2, with the
source's product-comparison proof retained separately from the earlier
logarithm proof.
\end{hblemma}

\begin{hbtheorem}{thm:p15-clt-iid}{proved=annex:p15.clt-iid}{The independent identically distributed central limit theorem}
If \(X_1,X_2,\ldots\) are i.i.d.,
\(\mathbb EX_1=\mu\), and
\(\operatorname{Var}(X_1)=\sigma^2\in(0,\infty)\), then, for
\(S_n=\sum_{j=1}^nX_j\),
\[
 \frac{S_n-n\mu}{\sigma\sqrt n}\Rightarrow\chi.
\]
Equivalently, \((S_n-n\mu)/\sqrt n\Rightarrow\sigma\chi\).
The latter statement also holds when \(\sigma=0\), in which case
every centered sum is zero almost surely. The ratio with \(\sigma\)
in its denominator is used only when \(\sigma>0\).
\end{hbtheorem}

\begin{hblemma}{thm:p15-clt-perturbation}{proved=annex:p15.clt-perturbation}{Perturbations and random denominators}
If \(Z_n\Rightarrow Z\) and \(R_n\to0\) in probability on the
same respective probability spaces, then \(Z_n+R_n\Rightarrow Z\).
If additionally \(A_n\to a\) and \(B_n\to b\) in probability,
where \(a,b\) are finite constants, then
\(A_nZ_n+B_n\Rightarrow aZ+b\).
If \(a\ne0\), then \(Z_n/A_n\Rightarrow Z/a\), assigning any
fixed finite value on \(\{A_n=0\}\).
No independence between these variables is assumed.
In particular a sequence of deterministic constants cannot converge
weakly to a nondegenerate normal law.
\end{hblemma}

\begin{hbtheorem}{thm:p15-clt-numerical-examples}{proved=annex:p15.clt-numerical-examples}{Roulette, binomial counts, dice, and Poisson counts}
The following calculations are for the stated independent trial models.
Finite normal approximations are estimates, not identities.
\begin{enumerate}
\item A roulette gain \(X\) with probabilities \(18/38\) at \(1\)
and \(20/38\) at \(-1\) has mean \(-1/19\) and variance
\(360/361\). After \(361\) plays, the source approximation to
\(\mathbb P(S_{361}\geq0)\) is
\(1-\Phi(19/\sqrt{360})\), about \(0.1583\);
replacing the standard deviation by one gives about \(0.1587\).
The expected gain is exactly \(-19\).
\item For fair coin indicators, \(S_n\) is binomial with parameters
\(n,1/2\), and
\((S_n-n/2)/(\sqrt n/2)\Rightarrow\chi\).
The source approximation for \(4900\leq S_{10000}\leq5100\) is
\(\Phi(2)-\Phi(-2)\), about \(0.9545\).
\item For integer-valued \(S\), histogram correction represents
\(\{S=k\}\) by the interval \([k-1/2,k+1/2]\) and
\(\{S\leq k\}\) by the threshold \(k+1/2\).
For \(S_{16}\) above, this gives
\[
 \mathbb P(S_{16}=8)\approx2\Phi(1/4)-1\approx0.1974,
 \qquad \mathbb P(S_{16}=8)=12870/65536.
 \]
For \(180\) fair die rolls the number \(D\) of sixes has mean
\(30\), variance \(25\), and
\[
 \mathbb P(D<25)\approx
 \Phi((24.5-30)/5)=\Phi(-1.1)\approx0.1357.
 \]
\item If \(Z_\lambda\) has Poisson law of mean \(\lambda>0\), then
\[
 (Z_\lambda-\lambda)/\sqrt\lambda\Rightarrow\chi
                         \quad(\lambda\to\infty\text{ through real values}).
 \]
\end{enumerate}
This supplies Examples 3.4.5--3.4.8 and Exercise 3.4.1.
\end{hbtheorem}

\begin{hbtheorem}{thm:p15-clt-pairwise-boundary}{proved=annex:p15.clt-pairwise-boundary}{Pairwise independence does not give the central limit theorem}
There is an identically distributed, pairwise independent sequence
\((X_j)\) of fair signs for which, with \(S_n=\sum_{j=1}^nX_j\),
for every integer \(r\geq0\),
\[
 \mathbb P(S_{2^r}=2^r)=\mathbb P(S_{2^r}=-2^r)=2^{-r-1},
 \qquad \mathbb P(S_{2^r}=0)=1-2^{-r}.
 \]
Thus \(S_{2^r}/\sqrt{2^r}\to0\) in probability although
\(\operatorname{Var}(S_n)=n\). In particular the i.i.d. normal
limit does not follow from pairwise independence.
The proof supplies the full recursive construction of Example 3.4.9.
\end{hbtheorem}

\begin{hbtheorem}{thm:p15-clt-oscillation}{proved=annex:p15.clt-oscillation}{Normal fluctuations do not settle to a probability limit}
For i.i.d. centered variables with variance
\(\sigma^2\in(0,\infty)\),
\[
 \limsup_{n\to\infty}S_n/\sqrt n=+\infty,\qquad
 \liminf_{n\to\infty}S_n/\sqrt n=-\infty
                                                \quad\text{almost surely}.
 \]
The sequence \(S_n/\sqrt n\) does not converge in probability
to any finite random variable. The proof follows Exercise 3.4.2's
tail-zero-one and factorial-subsequence routes.
\end{hbtheorem}

\begin{hbtheorem}{thm:p15-clt-sqrt-scale-necessity}{proved=annex:p15.clt-sqrt-scale-necessity}{A proper square-root-scale limit requires a finite second moment}
If \(X_j\) are i.i.d. and \(S_n/\sqrt n\) converges weakly to
a probability law on \(\mathbb R\), then
\(\mathbb EX_1^2<\infty\) and \(\mathbb EX_1=0\).
The limit is \(N(0,\mathbb EX_1^2)\), with the degenerate case
allowed. The proof includes Exercise 3.4.3's independent-copy
and symmetric truncation argument, including its conditional-sign step.
\end{hbtheorem}

\begin{hbtheorem}{thm:p15-clt-transformed-sums}{proved=annex:p15.clt-transformed-sums}{Square roots and self-normalized sums}
For i.i.d. \(X_j\geq0\), with mean \(1\) and variance
\(\sigma^2\in(0,\infty)\),
\[
 2(\sqrt{S_n}-\sqrt n)\Rightarrow\sigma\chi.
\]
For i.i.d. centered \(X_j\) with
\(\mathbb EX_j^2=\sigma^2\in(0,\infty)\), let
\[
 T_n=\frac{\sum_{j=1}^nX_j}{(\sum_{j=1}^nX_j^2)^{1/2}},
 \]
and set \(T_n=0\) if its denominator is zero. Then
\(T_n\Rightarrow\chi\).
These are Exercises 3.4.4--3.4.5 with all denominator conventions.
\end{hbtheorem}

\begin{hbtheorem}{thm:p15-clt-random-index}{proved=annex:p15.clt-random-index}{Random indices may depend on the observations}
Let \(X_j\) be i.i.d., centered, with variance
\(\sigma^2\in(0,\infty)\), and set \(S_0=0\).
Let \(N_n\) be finite nonnegative integer-valued random variables,
and let \(a_n\) be positive integers with \(a_n\to\infty\).
If \(N_n/a_n\to1\) in probability, then
\[
 S_{N_n}/(\sigma\sqrt{a_n})\Rightarrow\chi.
 \]
No stopping-time property and no independence of \(N_n\) from
the observations is required. The conclusion is valid also for
a continuously indexed family as its index tends to infinity,
by the same estimates.
\end{hbtheorem}

\begin{hbdefinition}{def:p15-clt-renewal}{Renewal notation for the fluctuation theorem}
Let \(Y_j\) be i.i.d. nonnegative waiting times with
\(\mu=\mathbb EY_1\in(0,\infty)\) and
\(\sigma^2=\operatorname{Var}(Y_1)\in(0,\infty)\).
Set \(T_0=0\), \(T_m=\sum_{j=1}^mY_j\), and
\(N(t)=\sup\{m\geq0:T_m\leq t\}\).
The renewal inverse law gives one event of probability one on
which \(T_m\to\infty\), \(N(t)<\infty\) for all \(t\geq0\),
and \(N(t)/t\to1/\mu\). Set \(N(t)=0\) on its null complement.
This convention counts all renewals in a group of tied epochs.
It includes the source's strictly positive waiting-time regime.
\end{hbdefinition}

\begin{hbtheorem}{thm:p15-clt-renewal-random-index}{proved=annex:p15.clt-renewal-random-index}{The renewal central limit theorem by random indexing}
With the preceding assumptions, as real \(t\to\infty\),
\[
 \frac{\mu N(t)-t}{\sigma\sqrt{t/\mu}}\Rightarrow\chi.
 \]
The proof uses Exercise 3.4.6's random-index theorem and proves
that the last incomplete waiting time is negligible on the
\(\sqrt t\) scale. This is the source route of Exercise 3.4.7.
\end{hbtheorem}

\begin{hbtheorem}{thm:p15-clt-renewal-interpolation}{proved=annex:p15.clt-renewal-interpolation}{The distinct deterministic-window renewal proof}
Under the same assumptions put
\(u=\lfloor t/\mu\rfloor\), \(D_t=T_u-t\), and
\(M_t=\lfloor t^{3/5}\rfloor\). Then
\[
 \mathbb P\left(
  \max_{\substack{m\in\mathbb Z,\ |m|\leq M_t\\u+m\geq0}}
  |T_{u+m}-T_u-m\mu|>t^{2/5}\right)\longrightarrow0,
 \qquad
 \frac{|N(t)-(t-D_t)/\mu|}{\sqrt t}\longrightarrow0
                                        \quad\text{in probability}.
 \]
These estimates yield the renewal central limit theorem again,
retaining the separate proof requested in Exercise 3.4.8.
\end{hbtheorem}

\section{Triangular arrays and their two normal-limit proofs}
\label{sec:p15-clt-arrays}

\begin{hbtheorem}{thm:p15-clt-lindeberg-feller}{proved=annex:p15.clt-lindeberg-feller}{Lindeberg--Feller for finite rows}
Let each row consist of independent centered square-integrable entries.
If \(v_n\to\sigma^2\in(0,\infty)\) and
\(L_n(\varepsilon)\to0\) for every \(\varepsilon>0\), then
\[
 S_n\Rightarrow\sigma\chi.
 \]
The theorem holds for arbitrary finite row lengths \(k_n\);
the source's \(k_n=n\) is included. Its conditions imply
\(\max_{j\leq k_n}v_{n,j}\to0\).
If instead \(v_n\to0\), then \(S_n\to0\) in mean square
without a separate Lindeberg assumption.
The annex first gives the source's characteristic-function proof.
\end{hbtheorem}

\begin{hblemma}{thm:p15-clt-diagonal-truncation}{proved=annex:p15.clt-diagonal-truncation}{The source's vanishing truncation threshold}
Under the hypotheses of the preceding nondegenerate theorem, there
exist positive deterministic \(\varepsilon_n\downarrow0\) with
\(L_n(\varepsilon_n)\to0\). Define
\[
 Y_{n,j}=X_{n,j}\mathbf1_{\{|X_{n,j}|\leq\varepsilon_n\}},
 \quad Z_{n,j}=Y_{n,j}-\mathbb EY_{n,j},\quad
 U_n=\sum_{j=1}^{k_n}Z_{n,j}.
 \]
Then \(|Z_{n,j}|\leq2\varepsilon_n\),
\(\mathbb E(S_n-U_n)^2\to0\), and
\(\sum_j\mathbb EZ_{n,j}^2\to\sigma^2\).
This includes Lemma 3.4.15 and all the truncation steps used
in the second source proof.
\end{hblemma}

\begin{hbtheorem}{thm:p15-clt-moment-array-proof}{proved=annex:p15.clt-moment-array-proof}{The independent method-of-moments proof}
For centered independent finite rows \(Z_{n,j}\) satisfying
\(|Z_{n,j}|\leq b_n\to0\) and
\(\sum_j\mathbb EZ_{n,j}^2\to\sigma^2\in[0,\infty)\), put
\(U_n=\sum_jZ_{n,j}\). For every positive integer \(r\),
\[
 \mathbb EU_n^r\longrightarrow
 \begin{cases}
 0,&r\text{ odd},\\
 \sigma^{2q}(2q)!/(2^qq!),&r=2q.
 \end{cases}
 \]
Consequently \(U_n\Rightarrow\sigma\chi\).
Together with the preceding truncation lemma, this proves the
full Lindeberg--Feller theorem again, without using its first proof
to establish these moment limits.
\end{hbtheorem}

\begin{hbtheorem}{thm:p15-clt-records-and-cycles}{proved=annex:p15.clt-records-and-cycles}{Normal fluctuations of records and permutation cycles}
Let \(B_j\) be independent Bernoulli variables with
\(\mathbb P(B_j=1)=1/j\), and put \(H_n=\sum_{j=1}^n1/j\).
Then, for \(n\geq2\),
\[
 \frac{\sum_{j=1}^nB_j-H_n}{\sqrt{\log n}}\Rightarrow\chi,
 \qquad
 \frac{\sum_{j=1}^nB_j-\log n}{\sqrt{\log n}}\Rightarrow\chi.
 \]
This is the limit for the number of cycles in a uniform permutation
of \(n\) objects. It is also the limit for the number of strict upper
records among \(n\) i.i.d. observations with a continuous CDF.
The record indicators are mutually independent with probabilities
\(1/j\); their full rank argument is supplied.
\end{hbtheorem}

\begin{hbtheorem}{thm:p15-clt-three-series-source}{proved=annex:p15.clt-three-series-source}{Discharging the source's three-series necessity route}
Let \(X_j\) be independent and suppose \(\sum_jX_j\) converges
almost surely to a finite limit. For every \(A>0\), set
\(Y_j=X_j\mathbf1_{\{|X_j|\leq A\}}\). Then
\[
 \sum_j\mathbb P(|X_j|>A)<\infty,\qquad
 \sum_j\mathbb EY_j\text{ converges},\qquad
 \sum_j\operatorname{Var}(Y_j)<\infty.
 \]
The full central-limit contradiction of Example 3.4.12 is now
unconditional. It supplies the formerly deferred Lindeberg--Feller
premise in the random-series chapter's separate source route.
\end{hbtheorem}

\begin{hbtheorem}{thm:p15-clt-infinite-variance-example}{proved=annex:p15.clt-infinite-variance-example}{A normal limit with infinite variance}
Let \(X_j\) be i.i.d. with symmetric law and
\(\mathbb P(|X_1|>x)=x^{-2}\) for every \(x\geq1\).
Then \(\mathbb EX_1=0\), \(\mathbb EX_1^2=\infty\), but
\[
 \frac{\sum_{j=1}^nX_j}{\sqrt{n\log n}}\Rightarrow\chi.
 \]
For all sufficiently large \(n\), the source truncation
\(c_n=\sqrt n\log\log n\) gives
\[
 \mathbb P(\exists j\leq n:|X_j|>c_n)
       \leq(\log\log n)^{-2},\qquad
 \mathbb E(X_1^2;|X_1|\leq c_n)=2\log c_n\sim\log n.
 \]
Both the source's upper/lower tail estimates and this exact
truncated moment calculation are supplied.
\end{hbtheorem}

\begin{hbtheorem}{thm:p15-clt-escaping-variance}{proved=annex:p15.clt-escaping-variance}{Weak convergence need not preserve variance}
Let \(X_1\) be a fair sign. For \(j\geq2\), let independent \(X_j\)
have probabilities \(1/(2j^2)\) at each of \(\pm j\), and
\((1-j^{-2})/2\) at each of \(\pm1\). Then
\[
 \operatorname{Var}(S_n)/n\to2,\qquad
 S_n/\sqrt n\Rightarrow\chi.
 \]
For every fixed \(\varepsilon>0\), the array \(X_j/\sqrt n\)
has Lindeberg tail sum tending to \(1\), rather than zero.
This is Exercise 3.4.9.
\end{hbtheorem}

\begin{hbtheorem}{thm:p15-clt-lyapunov}{proved=annex:p15.clt-lyapunov}{Lyapunov's criterion and the bounded special case}
Let \(W_{n,j}\), \(1\leq j\leq k_n\), be independent within rows,
with finite means and variances, and put
\[
 a_n^2=\sum_j\operatorname{Var}(W_{n,j}).
 \]
Assume \(0<a_n<\infty\) for all sufficiently large \(n\).
If some \(\delta>0\) satisfies
\[
 a_n^{-2-\delta}
       \sum_j\mathbb E|W_{n,j}-\mathbb EW_{n,j}|^{2+\delta}
                         \longrightarrow0,
 \]
then \(\sum_j(W_{n,j}-\mathbb EW_{n,j})/a_n\Rightarrow\chi\).
There is no additional requirement that \(a_n\to\infty\).
In particular:
\begin{enumerate}
\item Independent \(|X_j|\leq M<\infty\) with
 \(\sum_j\operatorname{Var}(X_j)=\infty\) satisfy the normal
 limit after centering and dividing by their cumulative standard deviation.
\item Independent centered \(X_j\) with variance one and
 \(\sup_j\mathbb E|X_j|^{2+\delta}<\infty\) satisfy
 \(S_n/\sqrt n\Rightarrow\chi\).
\end{enumerate}
This proves all of Exercises 3.4.10--3.4.12 and their finite-row extension.
\end{hbtheorem}

\begin{hbtheorem}{thm:p15-clt-sparse-jumps}{proved=annex:p15.clt-sparse-jumps}{The three regimes of sparse signed jumps}
Let \(\beta>0\), and let independent \(X_j\) have masses
\(1/(2j^\beta)\) at \(j\) and \(-j\), and mass \(1-j^{-\beta}\)
at zero. Then:
\begin{enumerate}
\item If \(\beta>1\), only finitely many \(X_j\) are nonzero
 almost surely, so \(S_n\) converges almost surely.
\item If \(0<\beta<1\),
 \[
  S_n/n^{(3-\beta)/2}\Rightarrow(3-\beta)^{-1/2}\chi.
 \]
\item If \(\beta=1\), then \(S_n/n\Rightarrow Q\), where
 \[
 \mathbb Ee^{itQ}
  =\exp\left\{-\int_0^1\frac{1-\cos(tx)}x\,dx\right\}.
 \]
 This formula defines a proper nonnormal probability law.
\end{enumerate}
This supplies all parameter cases in Exercise 3.4.13.
\end{hbtheorem}

\section{Prime divisors and the Erd\H{o}s--Kac theorem}
\label{sec:p15-clt-prime-divisors}

\begin{hblemma}{thm:p15-clt-prime-reciprocals}{proved=annex:p15.clt-prime-reciprocals}{The prime reciprocal sum needed by the source}
A prime is an integer \(p\geq2\) whose positive divisors are
exactly \(1,p\). In this section every sum indexed by \(p\)
ranges over primes. For real \(x\to\infty\),
\[
 \sum_{p\leq x}\frac1p=\log\log x+O(1),\qquad
 \sum_p p^{-2}<\infty.
 \]
Here \(O(1)\) means a quantity bounded in absolute value for all
sufficiently large \(x\). The annex proves the reciprocal estimate
elementarily, including the needed bound on
\(\vartheta(x)=\sum_{p\leq x}\log p\).
It supplies the actual premise of Durrett's Erd\H{o}s--Kac proof;
it is not a claim to reproduce Hardy and Wright's separately
referred chapter.
\end{hblemma}

\begin{hbdefinition}{def:p15-clt-prime-count}{Uniform integers and distinct prime divisors}
For each positive integer \(n\), let \(\mathbb P_n\) be the
uniform law on \(\{1,\ldots,n\}\), with expectation \(\mathbb E_n\).
For a positive integer \(m\), write
\[
 \delta_p(m)=\mathbf1_{\{p\mid m\}},\qquad
 g(m)=\sum_{p\mid m}1,\qquad g(1)=0.
 \]
Thus \(g\) counts distinct prime divisors, without multiplicity.
For \(A\subseteq\mathbb N\), its natural density, when it exists,
is \(\lim_n\mathbb P_n(A\cap\{1,\ldots,n\})\).
These limits do not define a countably additive probability
measure on all subsets of \(\mathbb N\).
\end{hbdefinition}

\begin{hbtheorem}{thm:p15-clt-erdos-kac}{proved=annex:p15.clt-erdos-kac}{The source's prime-divisor central limit theorem}
For every \(x\in\mathbb R\),
\[
 \mathbb P_n\left(
   \frac{g(m)-\log\log n}{\sqrt{\log\log n}}\leq x\right)
                      \longrightarrow\Phi(x).
 \]
The formula is used for all sufficiently large \(n\), so its
denominator is real and positive. The proof retains the source's
cutoff \(\alpha_n=n^{1/\log\log n}\), its comparison with
independent Bernoulli prime indicators, and its transfer of
every fixed centered moment. The actual divisibility indicators
under \(\mathbb P_n\) are not assumed independent.
\end{hbtheorem}

\section{A uniform quantitative normal approximation}
\label{sec:p15-clt-berry-esseen}

\begin{hbdefinition}{def:p15-clt-smoothing}{The source's smoothing distribution}
For \(L>0\) put
\[
 h_L(x)=\frac{1-\cos(Lx)}{\pi Lx^2}\quad(x\ne0),
 \qquad h_L(0)=L/(2\pi).
 \]
Let \(H_L\) be its CDF. This is the scaled Pólya density;
its characteristic function is
\(\omega_L(t)=(1-|t|/L)^+\).
For a bounded Borel function \(D\), define
\((D*H_L)(x)=\int D(x-y)h_L(y)\,dy\).
For a CDF \(F\), this is the CDF of the sum of independent
variables with CDFs \(F,H_L\).
For any two CDFs define their uniform distance by
\(d(F,G)=\sup_{x\in\mathbb R}|F(x)-G(x)|\).
\end{hbdefinition}

\begin{hblemma}{thm:p15-clt-smoothing-bound}{proved=annex:p15.clt-smoothing-bound}{Smoothing does not hide an arbitrary CDF discrepancy}
Suppose \(F,G\) are CDFs and
\(|G(y)-G(x)|\leq\lambda|y-x|\) for a finite
\(\lambda>0\). Let
\(\eta=d(F,G)\), \(D=F-G\), and
\(\eta_L=\sup_x|(D*H_L)(x)|\). Then
\[
 \eta_L\geq\eta/2-\frac{12\lambda}{\pi L},
 \qquad
 \eta\leq2\eta_L+\frac{24\lambda}{\pi L}.
 \]
In particular the lemma applies to a differentiable CDF \(G\)
with \(0\leq G'\leq\lambda\).
\end{hblemma}

\begin{hblemma}{thm:p15-clt-cdf-inversion}{proved=annex:p15.clt-cdf-inversion}{CDF inversion and the smoothing inequality}
Let \(K_1,K_2\) be CDFs whose characteristic functions
\(\kappa_1,\kappa_2\) are integrable on \(\mathbb R\).
If \((\kappa_1(t)-\kappa_2(t))/t\) is integrable as well
(its value at \(t=0\) is immaterial), then
\[
 K_1(x)-K_2(x)=\frac1{2\pi}
       \int_{\mathbb R}-e^{-itx}
                 \frac{\kappa_1(t)-\kappa_2(t)}{it}\,dt.
 \]
Finite first absolute moments for both laws suffice for this
quotient condition; in particular the source's integrable,
mean-zero assumption suffices.
If \(F,G\) have finite first absolute moments and \(G\) obeys
the preceding Lipschitz bound, then for every \(L>0\),
\[
 d(F,G)\leq\frac1\pi\int_{-L}^L
       \frac{|\varphi_F(t)-\varphi_G(t)|}{|t|}\,dt
                      +\frac{24\lambda}{\pi L}.
 \]
The quotient condition is proved directly for the smoothed
transforms; the convolved laws need not have first moments.
\end{hblemma}

\begin{hbtheorem}{thm:p15-clt-berry-esseen}{proved=annex:p15.clt-berry-esseen}{The source's Berry--Esseen bound with constant three}
Let \(X_j\) be i.i.d., centered, with
\(\mathbb EX_1^2=\sigma^2\in(0,\infty)\) and
\(\rho=\mathbb E|X_1|^3<\infty\).
If \(F_n\) is the CDF of \(\sum_{j=1}^nX_j/(\sigma\sqrt n)\),
then, for every integer \(n\geq1\),
\[
 \sup_{x\in\mathbb R}|F_n(x)-\Phi(x)|
                         \leq\frac{3\rho}{\sigma^3\sqrt n}.
 \]
The proof retains the source's smoothing kernel, cutoff,
product comparison, and explicit constant calculation.
The order \(n^{-1/2}\) cannot be improved uniformly even for
bounded fair-sign summands.
\end{hbtheorem}

\begin{proofboundary}{def:p15-clt-current-boundary}{The completed normal-domain result and historical continuations}
The results above include the source's i.i.d. and general finite-row
central limit theorems, both array proof routes, Examples 3.4.5--3.4.13,
Exercises 3.4.1--3.4.13, Erd\H{o}s--Kac, and the full
Berry--Esseen proof with Lemmas 3.4.18--3.4.19.
The normal-domain criterion, Theorem~3.4.14, is proved in
\cref{thm:p15-nda-characterization}. For an i.i.d. law that is nondegenerate,
meaning not concentrated at a single point, it
characterizes the existence of deterministic \(a_n\in\mathbb R\),
\(b_n>0\), with \((S_n-a_n)/b_n\Rightarrow\chi\), by
\[
 \frac{y^2\mathbb P(|X_1|>y)}
      {\mathbb E(X_1^2;|X_1|\leq y)}\longrightarrow0.
 \]
Its denominator is used only for sufficiently large \(y\) where
it is positive. Both directions and all admissible center/scale
choices are supplied there; \cref{thm:p15-nda-normalizer-construction}
constructs the normalizers. The exact 1954 Section~35 cutoff route is
\cref{thm:p15-nda-gk-cutoff-route}, with its complete i.i.d. interface.
The broader original general-array and canonical-measure presentations
are distinguished in \cref{def:p15-nda-historical-boundary}.
The additional rate equivalence cited in the source's
Heyde remark, its separate historical proofs, and the later
local-limit and stable-law sections remain separate presentations.
\end{proofboundary}
